\section{位置必须外挂，因为注意力看不见顺序}
\label{sec:14-Deep-Learning/05-Attention-and-Sequence-Models/03}

你从头写了一个小 Transformer，任务简单到不该出错：给定一串数，判断它是不是升序。数据是自己生成的，标签绝对干净，模型能过拟合训练集里的单个 batch，但验证准确率死死钉在 \(50\%\)。你怀疑数据泄漏、怀疑学习率、怀疑初始化，最后随手做了一件事——把同一条序列的元素打乱再喂进去，发现输出的 logit 一个比特都没变。

不是训练出了问题。模型根本没有能力知道谁在前谁在后。这件事可以精确证明。

\subsection{置换等变是注意力的结构性质}

把不带位置信息的单头注意力写成一个映射 \(\operatorname{Attn}:\R^{n\times d}\to\R^{n\times d}\)：

\[
\operatorname{Attn}(X) = \softmax_{\text{行}}\!\left(\frac{XW_QW_K^{\trans}X^{\trans}}{\sqrt{d_k}}\right) XW_V .
\]

\begin{proposition}
设 \(P\in\R^{n\times n}\) 是任意置换矩阵。若不加位置编码且不加任何位置相关的 mask，则
\[
\operatorname{Attn}(PX) = P\,\operatorname{Attn}(X).
\]
\end{proposition}

证明。三组投影都是右乘常数矩阵，与左乘 \(P\) 可交换，所以 \(Q\mapsto PQ\)、\(K\mapsto PK\)、\(V\mapsto PV\)。分数矩阵

\[
\begin{aligned}
S' &= \frac{(PQ)(PK)^{\trans}}{\sqrt{d_k}}
   = \frac{PQK^{\trans}P^{\trans}}{\sqrt{d_k}}
   = P S P^{\trans}.
\end{aligned}
\]

设 \(P\) 对应下标置换 \(\pi\)，则 \((PSP^{\trans})_{ij}=S_{\pi(i)\pi(j)}\)：\(S'\) 的第 \(i\) 行不过是 \(S\) 的第 \(\pi(i)\) 行把元素重新排了个序。行 softmax 逐行作用，且对分量重排是等变的（分母是同一批数的和，重排不改变和），于是 \(A'=PAP^{\trans}\)。最后

\[
A'V' = (PAP^{\trans})(PV) = PA(P^{\trans}P)V = PAV = P\operatorname{Attn}(X),
\]

用到 \(P^{\trans}P=I\)。证毕。

条件对账。这个证明用到三件事，每一件都指向一种破坏对称的办法。\(W_Q,W_K,W_V\) 在所有位置共享——权重共享正是等变性的来源，如果每个位置用不同的投影矩阵，命题立刻失效（早期的一些位置相关线性层就是这个思路，随之而来的是无法处理变长序列）。softmax 沿行作用且对分量重排等变——这一步无可绕开。mask 与位置无关——因果 mask 恰恰违反这一条，它是唯一``免费''的位置信息来源。

推论比命题更值得记住。逐位置的前馈层、LayerNorm、残差连接全都是逐行作用的，同样满足等变性，而等变映射的复合仍然等变。所以\emph{整个} Transformer 编码器，无论堆多少层，都是置换等变的。若读出头是置换不变的（比如平均池化），那么``猫追狗''和``狗追猫''会得到逐比特相同的预测——不是相近，是相同。

这就是为什么位置编码不是一个可选的工程改良。它是打破一个精确对称性的必需品。

\subsection{正弦编码：为什么频率排成几何级数}

原始 Transformer 给位置 \(m\) 分配一个固定向量，第 \(i\) 个二维块写成

\[
\bigl(\sin(\omega_i m),\ \cos(\omega_i m)\bigr),
\qquad
\omega_i = 10000^{-2i/d},\quad i=0,1,\ldots,\tfrac{d}{2}-1,
\]

然后把这个向量加到词嵌入上。频率 \(\omega_i\) 从 1 按几何级数递减到 \(10^{-4}\)。

几何级数不是随手挑的。固定一个位移 \(k\)，用和角公式：

\[
\begin{bmatrix}\sin(\omega_i(m+k))\\ \cos(\omega_i(m+k))\end{bmatrix}
=
\begin{bmatrix}\cos(\omega_i k) & \sin(\omega_i k)\\ -\sin(\omega_i k) & \cos(\omega_i k)\end{bmatrix}
\begin{bmatrix}\sin(\omega_i m)\\ \cos(\omega_i m)\end{bmatrix}.
\]

右边那个 \(2\times 2\) 矩阵只依赖 \(k\)，与 \(m\) 无关。把所有二维块摞起来，``向后移动 \(k\) 个位置''对应的是一个固定的分块旋转矩阵，同一个 \(k\) 在序列任何地方都对应同一个线性变换。这正是模型能用线性投影读出相对位置的原因——注意力里的 Query 与 Key 投影都是线性的，它们有机会学到近似这种旋转的算子。

频率的跨度决定了能分辨的尺度范围。最高频 \(\omega_0=1\) 的周期是 \(2\pi\approx 6.3\)，相邻位置之间就有明显相位差，负责精细区分；最低频 \(\omega_{d/2-1}\approx10^{-4}\) 的周期约 \(6.3\times10^4\)，在整个训练长度内都不绕回一圈，提供一个不重复的粗尺度坐标。几何级数让每个数量级分到大致相同数量的维度，是对数刻度上的均匀采样，与 \hyperref[sec:03-Linear-Algebra/08-Complex-Matrices-and-Fourier-Structure/03]{``Fourier 基与 FFT''}里用不同频率基分解信号是同一种想法。

\begin{center}
\begin{tikzpicture}[x=0.135cm, y=0.52cm]
% 取 d=24；TikZ 的 sin 以角度为单位，故把弧度频率乘 180/pi。
\foreach \j/\w/\lab in {0/57.296/{\(\omega_0\)}, 1/26.594/{\(\omega_1\)}, 2/12.344/{\(\omega_2\)},
                        3/5.730/{\(\omega_3\)}, 4/2.659/{\(\omega_4\)}, 5/1.234/{\(\omega_5\)}}{
  \foreach \p in {0,...,47}{
    \pgfmathtruncatemacro{\v}{round(50+47*sin(\p*\w))}
    \fill[black!\v] (\p,-\j) rectangle (\p+1,{-\j+0.82});
  }
  \node[anchor=east] at (-1.5,{-\j+0.41}) {\footnotesize \lab};
}
\draw[black!70] (0,0.82) rectangle (48,-5);
% 位移标记
\draw[dashed,black!80] (12,1.5) -- (12,-5.6);
\draw[dashed,black!80] (20,1.5) -- (20,-5.6);
\draw[->,>=stealth] (12,1.2) -- (20,1.2);
\node[anchor=south] at (16,1.3) {\footnotesize 位移 \(k\)};
\node[anchor=north] at (12,-5.7) {\footnotesize 位置 \(m\)};
\node[anchor=west] at (49,0.4) {\footnotesize 高频：相邻位置即可分辨};
\node[anchor=west] at (49,-4.6) {\footnotesize 低频：整段序列内不绕回};
\end{tikzpicture}
\end{center}

\emph{图取 \(d=24\) 时前六个频率块，深浅代表正弦编码随位置的取值。最高频的周期确为 \(2\pi\approx6.3\) 个位置；频率随后按 \(10000^{-2/d}\) 的公比递减。两条虚线之间的位移 \(k\) 在每一个频率块上都对应同一个旋转角 \(\omega_i k\)，与 \(m\) 在哪里无关——相对位置因此可以被一个固定的线性变换读出。}

最下面那几行看起来只是一条平缓的明暗渐变，没有条纹，这不是画错了：低频维度在 48 个位置的窗口里连一个周期都走不完。它们在近距离上几乎不提供分辨力，作用是在跨越几千个位置时仍给出一条单调可辨的粗坐标。最上面几行正相反，相邻位置就有明显反差，但走出几十个位置之后相位早已绕回多圈，单独看无法判断绝对位置。两端都不够用，几何级数把中间的所有尺度都填上，才让位置既能精细区分又不混淆。

\subsection{可学习编码丢掉的正是外推}

另一种做法是干脆学一张表：开一个 \(n_{\max}\times d\) 的参数矩阵，第 \(m\) 行就是位置 \(m\) 的编码，跟着梯度一起训。在训练长度以内，它的表现通常不比正弦编码差，甚至略好——它可以拟合任务真正需要的位置结构，而不必迁就三角函数的形状。

短板在长度一超出 \(n_{\max}\) 时就暴露：第 \(n_{\max}+1\) 行不存在。这不是``效果变差''，是根本没有定义。正弦公式对任意实数 \(m\) 都能算出向量，可学习表则只有那 \(n_{\max}\) 行。工程上的补救是位置插值——把新长度上的位置线性映射回 \([0,n_{\max}]\) 再查表并插值，但这改变了相邻位置之间的间距，通常需要额外微调才能收敛。

正弦编码能算出任意位置的向量，也不等于模型知道怎么用它。训练时从未见过 \(m=8000\) 的相位组合，Query 与 Key 投影就没有在那一带被约束过，外推质量照样下降。这是一条经验规律，不是定理：能算出编码是必要条件，不是充分条件。

\subsection{RoPE 把相对位置写进旋转}

正弦编码和可学习编码都是\emph{加}到输入上的，位置信息与内容信息挤在同一个加法通道里，中间还要穿过 \(W_Q,W_K\) 才影响到内积。RoPE 换了个作用点：不动输入，直接对 Query 和 Key 施加位置相关的旋转。

把 \(q,k\in\R^{d_k}\) 按相邻两维分成 \(d_k/2\) 个二维块，第 \(i\) 块在位置 \(m\) 上转过角度 \(m\theta_i\)（\(\theta_i\) 仍取几何级数）。记这个分块旋转为 \(R_m\)，则位置 \(m\) 的查询与位置 \(m'\) 的键的内积是

\[
(R_m q)^{\trans}(R_{m'} k) = q^{\trans}R_m^{\trans}R_{m'} k = q^{\trans}R_{\,m'-m}\,k .
\]

关键一步是 \(R_m^{\trans}R_{m'}=R_{m'-m}\)：同一平面内的旋转关于角度相加构成一个交换群，转置就是反向旋转。分数因此\emph{只}依赖位置差 \(m'-m\)，绝对位置在内积里被彻底消掉。

这解决了加法编码的一个结构隐患：加法编码里，位置向量和词向量共用同一批坐标，模型必须自己学会把两者分开；RoPE 则把位置作用在匹配这一步，Value 完全不受影响——读取的内容始终是纯粹的内容。它也天然地是相对的，不需要模型再去学``相对位置怎么从两个绝对位置里减出来''。

边界同样要说清。RoPE 在超出训练长度时不会像可学习表那样直接失效，但远距离的旋转角 \(m\theta_i\) 会跑到训练中从未出现过的区间，注意力分数的分布随之漂移，长度外推仍然需要缩放 \(\theta_i\) 一类的额外处理。位置写进旋转，只是把问题从``有没有定义''换成了``在不在见过的范围内''，没有把它取消。

三种做法共享一个前提：注意力自身对顺序无知，位置必须从外面补。下一篇处理另一件必须补上的事——就算位置信息给足了，把这样的层堆到几十上百层，梯度还是会在中途散掉。
